\chapter{Belajar dari Jejak Pengguna}
\label{ch:rekomendasi}

Kuliah learning from user-generated data saya ikuti pada musim panas 2025, diberikan oleh kelompok
Markus Schedl, yang banyak meneliti rekomendasi musik. Sebelas kuliahnya bergerak dari collaborative
filtering klasik ke evaluasi, data teks dan data implisit, sistem rekomendasi berbasis graf, hibrida,
dan konteks, lalu ditutup dengan dua topik yang jarang dibahas di kuliah teknis: bias dan keadilan, serta
sistem rekomendasi yang diilhami psikologi. Latihannya membangun rekomendasi item-kNN dari umpan balik
implisit, model faktor laten, dan evaluasi dengan nDCG.

Sistem rekomendasi barangkali adalah machine learning yang paling sering kita temui tanpa disadari:
video berikutnya, lagu berikutnya, produk yang ``mungkin juga Anda suka''. Datanya adalah jejak yang
ditinggalkan jutaan pengguna, berupa rating, klik, putaran lagu, dan waktu menonton, dan tugasnya adalah
menebak apa yang akan disukai seseorang dari apa yang disukai orang-orang lain.

\section{Collaborative filtering}

Data dasarnya adalah matriks pengguna terhadap item, dengan rating di sel yang diketahui dan kosong di
sebagian besar sel lainnya. \istilah{Collaborative filtering} berbasis pengguna mengikuti logika sederhana:
cari pengguna lain yang selera masa lalunya mirip denganmu, lalu rekomendasikan apa yang mereka sukai.
Kemiripan biasanya diukur dengan kosinus atau korelasi Pearson antara vektor rating. Untuk dua pengguna
dengan rating $(5,3,4)$ dan $(4,2,5)$ pada tiga item yang sama, kemiripan kosinusnya
\begin{equation}
\frac{5\cdot4+3\cdot2+4\cdot5}{\sqrt{5^2+3^2+4^2}\sqrt{4^2+2^2+5^2}}=\frac{46}{\sqrt{50}\sqrt{45}}\approx0{,}97.
\end{equation}
Korelasi Pearson mengurangkan rata-rata rating setiap pengguna lebih dulu, sehingga pengguna yang selalu
memberi rating tinggi dan pengguna yang pelit rating tetap bisa dikenali sebagai bersepakat. Prediksi
rating lalu dihitung sebagai rata-rata berbobot dari rating tetangga-tetangga terdekat.

Versi berbasis item \citep{sarwar2001} membalik sudut pandangnya: item yang sering disukai oleh orang
yang sama dianggap mirip, dan kita merekomendasikan item yang mirip dengan yang sudah kamu sukai.
Karena jumlah item biasanya lebih stabil daripada jumlah pengguna, kemiripan antar item bisa dihitung
lebih jarang dan dipakai ulang, sehingga pendekatan ini lebih mudah diskalakan.

\section{Faktor laten}

Pendekatan yang menjadi inti solusi-solusi terbaik dalam kompetisi Netflix Prize adalah faktorisasi
matriks \citep{koren2009}. Setiap
pengguna $u$ dan setiap item $i$ diwakili oleh vektor berdimensi kecil, $\bm p_u$ dan $\bm q_i$, dan rating
diprediksi sebagai
\begin{equation}
\hat r_{ui}=\mu+b_u+b_i+\bm p_u^\top\bm q_i,
\end{equation}
dengan $\mu$ rata-rata global, $b_u$ kecenderungan pengguna memberi rating tinggi atau rendah, dan $b_i$
kecenderungan item dinilai tinggi atau rendah. Vektor-vektor ini dilatih dengan meminimalkan galat
kuadrat pada rating yang diketahui saja, ditambah regularisasi L2:
\begin{equation}
\Loss=\sum_{(u,i)\ \text{diketahui}}\big(r_{ui}-\hat r_{ui}\big)^2+\lambda\big(b_u^2+b_i^2+\|\bm p_u\|^2+\|\bm q_i\|^2\big).
\end{equation}
Turunan terhadap $\bm p_u$ untuk satu rating, dengan galat $e_{ui}=r_{ui}-\hat r_{ui}$, adalah
$-2e_{ui}\bm q_i+2\lambda\bm p_u$. Langkah SGD (\cref{ch:optimasi}) dengan laju $\eta$, faktor dua diserap ke
dalam $\eta$, menjadi
\begin{equation}
\bm p_u\leftarrow\bm p_u+\eta\big(e_{ui}\bm q_i-\lambda\bm p_u\big),\qquad
\bm q_i\leftarrow\bm q_i+\eta\big(e_{ui}\bm p_u-\lambda\bm q_i\big).
\end{equation}
Dimensi-dimensi laten tidak diberi nama, tetapi sering bisa ditafsirkan: untuk film, satu dimensi
mungkin memisahkan film serius dari komedi, dimensi lain film untuk penonton dewasa dari film keluarga.
Ini analisis faktor dari \cref{ch:unsup}, diterapkan pada matriks yang sebagian besar kosong.

Dalam praktik, kebanyakan data bukan rating, melainkan umpan balik implisit: klik, putaran, pembelian.
Data ini tidak punya sinyal negatif yang jelas, karena tidak mengklik bisa berarti tidak suka atau
sekadar tidak melihat. Bayesian personalized ranking \citep{rendle2009} mengubah tujuannya menjadi
peringkat: untuk setiap pengguna, item yang pernah ia pakai seharusnya diberi skor lebih tinggi dari item
yang belum. Loss-nya $-\sum\ln\sigma(\hat x_{ui}-\hat x_{uj})$ untuk pasangan item $i$ yang dipakai dan $j$
yang tidak, sebuah binary cross-entropy atas pasangan. Neural collaborative filtering \citep{he2017}
mengganti hasil kali titik dengan jaringan saraf, dan pendekatan berbasis graf seperti LightGCN
\citep{he2020} memandang pengguna dan item sebagai simpul graf bipartit dan menyebarkan embedding lewat
message passing (\cref{ch:gdl}).

\section{Mengukur rekomendasi yang baik}

Galat rating seperti RMSE kurang relevan kalau yang ditampilkan adalah daftar sepuluh teratas. Ukuran yang
peka peringkat lebih cocok. Precision dan recall pada $k$ teratas menghitung berapa item relevan yang
masuk daftar. \istilah{Discounted cumulative gain} \citep{jarvelin2002} menambahkan bahwa item relevan di
urutan atas lebih berharga daripada di urutan bawah,
\begin{equation}
\mathrm{DCG}@k=\sum_{i=1}^k\frac{\mathrm{rel}_i}{\log_2(i+1)},
\end{equation}
dan versi ternormalisasinya, nDCG, membaginya dengan DCG dari urutan terbaik yang mungkin. Untuk daftar
dengan relevansi $(3,2,0,1)$, DCG-nya $3/1+2/1{,}585+0/2+1/2{,}322\approx3+1{,}262+0+0{,}431=4{,}693$. Urutan
idealnya $(3,2,1,0)$ dengan DCG $3+1{,}262+0{,}5+0=4{,}762$. Maka nDCG $\approx0{,}985$: hampir sempurna,
hanya item relevan terakhir yang sedikit terlambat.

Kuliah menekankan bahwa akurasi bukan segalanya. Daftar yang akurat tetapi hanya berisi item yang sudah
populer tidak membantu siapa pun menemukan sesuatu yang baru. Ukuran di luar akurasi menilai keragaman
daftar, kebaruan, cakupan katalog, dan \istilah{serendipity}, yaitu kejutan yang menyenangkan. Evaluasi
offline pada data historis juga berbeda dari evaluasi online dengan pengguna sungguhan, dan studi dengan
pengguna sering memperlihatkan bahwa yang dirasakan pengguna tidak selalu sejalan dengan angka offline.

\section{Konten, konteks, dan hibrida}

Collaborative filtering gagal untuk item baru yang belum pernah dipakai siapa pun dan pengguna baru yang
belum punya riwayat, masalah yang disebut \istilah{cold start}. Rekomendasi berbasis konten mengisi celah ini
dengan memakai sifat item itu sendiri: teks deskripsi, tag, atau sinyal audio sebuah lagu. Untuk teks,
bobot TF-IDF memberi nilai tinggi pada kata yang sering muncul di sebuah dokumen tetapi jarang di
dokumen lain, dan model yang lebih maju memakai embedding kata dan model topik. Sistem hibrida
menggabungkan kedua sumber, misalnya dengan merata-ratakan skor, beralih di antara keduanya tergantung
situasi, atau memasukkan fitur konten ke dalam model faktor.

Konteks juga mengubah apa yang pantas direkomendasikan. Lagu yang cocok untuk berolahraga pagi berbeda
dari lagu untuk menjelang tidur. Kuliah memberi contoh dari riset kelompoknya sendiri tentang
rekomendasi musik yang memperhitungkan lokasi dan aktivitas, dengan memetakan tempat dan lagu ke tag
yang sama. Rekomendasi berbasis graf pengetahuan menyebarkan preferensi pengguna lewat relasi antar
item, seperti penyanyi, genre, dan label rekaman.

\section{Bias, keadilan, dan psikologi}

Sistem rekomendasi tidak hanya mencerminkan selera, tetapi juga membentuknya. \istilah{Popularity bias}
membuat item yang sudah populer semakin sering direkomendasikan, sehingga semakin populer lagi. Kinerja
sistem juga bisa berbeda antar kelompok pengguna, misalnya menurut usia, jenis kelamin, atau negara.
Kuliah membahas ukuran untuk mendeteksinya, seperti selisih kinerja rekomendasi antar kelompok, dan contoh
dari data musik tentang bagaimana algoritme rekomendasi bisa menggeser porsi musik lokal yang diterima
pendengar di berbagai negara. Strategi perbaikannya bisa diterapkan sebelum pelatihan, misalnya
menyeimbangkan ulang data, selama pelatihan dengan kendala keadilan, atau sesudahnya dengan menyusun
ulang daftar.

Kuliah terakhir menghubungkan rekomendasi dengan psikologi kognitif. Ingatan manusia memudar menurut
kurva lupa Ebbinghaus, dan lagu yang terakhir didengar beberapa waktu lalu lebih mudah diingat
kembali kalau sering diulang. Model memori dari arsitektur kognitif seperti ACT-R dipakai untuk
meramalkan lagu mana yang ingin didengar lagi seseorang, dengan menimbang seberapa sering dan seberapa
baru ia mendengarkannya. Sistem rekomendasi yang diilhami psikologi ini juga lebih mudah dijelaskan,
karena alasannya bisa diungkapkan dengan bahasa yang dipahami manusia.

\section{Di luar layar}

Rekomendasi adalah masalah keputusan berurutan, sehingga reinforcement learning (\cref{ch:rl}) mulai
dipakai untuk mengoptimalkan kepuasan jangka panjang, bukan klik berikutnya. Gagasan faktorisasi matriks
juga muncul di bidang lain: matriks senyawa terhadap target di penemuan obat (\cref{ch:hayati}) sama
jarangnya dengan matriks pengguna terhadap film, dan model yang belajar dari banyak target sekaligus
memakai logika yang sama dengan collaborative filtering. Dan pertanyaan tentang bias serta dampak
sosial rekomendasi membawa kita ke bab berikutnya.

\sumberbab{Bab ini mengikuti kuliah dan latihan Learning from User-generated Data: collaborative filtering,
faktor laten, evaluasi, data konten dan implisit, rekomendasi berbasis graf, hibrida, konteks, bias dan
keadilan, serta sistem rekomendasi yang diilhami psikologi.}
